\chapter{Математическое моделирование}\label{ch:model}
Производственный план должен не только экономно использовать материал,
но и сохранять возможность исполнения при отклонениях длительности операций и временной недоступности оборудования. Даже допустимое по расчётным временам расписание может приводить к задержкам заказов, накоплению полуфабриката и простоям, если фактический ход производства отличается от расчётного. Поэтому качество планирования рассматривается на двух уровнях: исходная эффективность определяет перспективные варианты организации производства,
а имитационное моделирование позволяет оценивать их устойчивость к заданным возмущениям.

Под устойчивостью здесь понимается способность фиксированного плана завершаться без блокировок и выполнять заказы в срок в рассматриваемых условиях исполнения,
 а также ограничивать величину задержек. Это операционная характеристика плана относительно выбранной модели случайных факторов.
 В главе последовательно изложены общая схема исследования, исходная математическая основа и дискретно-событийная модель, связывающая отдельные прогоны со статистическими оценками и выбором плана.

\section{Общая схема формирования и оценки планов}\label{sec:model/workflow}
Численные методы сначала формируют несколько различных допустимых производственных планов по расчётным характеристикам оборудования. Поиск и оценка устойчивости разделены: после завершения поиска каждый сохранённый вариант многократно исполняется в имитационной модели при одинаковых для сравниваемых вариантов сценарных предположениях. Затем рассчитываются показатели завершения производства, соблюдения сроков, загрузки оборудования и складов. Окончательный выбор производится по правилу, зафиксированному до проведения опытов, с учётом исходного качества и неопределённости статистических оценок.

\begin{center}
\begin{tikzpicture}[
  modelstage/.style={draw, rounded corners, align=center, text width=10.4cm, minimum height=9mm},
  modellink/.style={-{Latex}, thick}, node distance=5mm, font=\small]
  \node[modelstage] (methods) {Численные методы: поиск по исходным критериям};
  \node[modelstage, below=of methods] (plans) {Несколько различных допустимых производственных планов};
  \node[modelstage, below=of plans] (runs) {Имитация исполнения каждого плана в заданных сценариях};
  \node[modelstage, below=of runs] (metrics) {Оценки сроков, задержек, блокировок, загрузки и их неопределённости};
  \node[modelstage, below=of metrics] (choice) {Выбор плана по заранее установленному правилу};
  \draw[modellink] (methods) -- (plans);
  \draw[modellink] (plans) -- (runs);
  \draw[modellink] (runs) -- (metrics);
  \draw[modellink] (metrics) -- (choice);
\end{tikzpicture}
\captionof{figure}{Общая схема формирования и имитационной оценки производственных планов}
\label{fig:model/workflow}
\end{center}

Такая последовательность не подменяет оптимизацию случайным проигрыванием производства. Численные методы отвечают за построение содержательных альтернатив, а имитация выявляет различия между ними, не видимые при сравнении только расчётных показателей. Небольшая выборка удачных прогонов не доказывает устойчивость плана: вместе с оценками должны указываться число повторений, частота незавершения и статистическая неопределённость. Если приемлемых вариантов не найдено, это фиксируется как результат оценки, а не маскируется выбором формально лучшего из неприемлемых планов.

Исходная модель необходима для задания состава операций, материальных связей и ограничений, которые имитация должна сохранять при изменении фактических времён исполнения.
Настоящая работа предлагает модель, в которой одновременно учитываются: требования по срокам, ограничения по складам, по переналадкам ножей и очередности схем раскроя --- при одновременной минимизации нескольких критериев эффективности.
Дополнительно постановка учитывает требование размещения взаимодействующего оборудования (БДМ, ПРС/БРС и складов) в одном производственном помещении, а устойчивость получаемых планов к случайным возмущениям оценивается с помощью разработанной имитационной модели исполнения расписания.
В известных публикациях такая комплексная постановка ранее не рассматривалась.
% Тем самым обзор литературы обосновывает выбор направления исследования и показывает место настоящей работы среди существующих подходов. С одной стороны, работа опирается на развитый аппарат классической задачи раскроя, методов целочисленного программирования и метаэвристик. С другой стороны, предлагается постановка, ориентированная именно на интеграцию раскроя, выпуска полуфабриката, складирования и календарного планирования в условиях целлюлозно-бумажного производства. Это создаёт основу для последующего изложения математической модели в главе~\ref{ch:model} и численных методов её решения в главе~\ref{ch:algorithm}.

\section{Постановка задачи}\label{sec:model/statement}
В данном разделе строится формальная математическая модель задачи. Она задаёт исходный план при расчётных длительностях операций, без случайных отказов и ремонтов; в имитационной модели технологические и складские ограничения сохраняются, а соблюдение сроков проверяется по фактическому исполнению.
Последовательно вводятся входные параметры, описывающие заказы и оборудование, переменные решения (планы выпуска, нарезки и расписание), а затем фиксируются целевая функция и ограничения выполнимости.

Модель охватывает полный производственный цикл: от выпуска тамбуров на бумагоделательных машинах, через складирование полуфабриката, до нарезки на продольно-резательных станках.
Такое построение позволяет учитывать взаимосвязи между этапами производства и формулировать задачу оптимизации как единое целое.


Входные параметры модели описывают характеристики заказов, оборудования и технологические константы.

\[
B = \{ b_{1},\ \ldots,\ b_{U}\}
\]
--- множество времен установок марок бумаги, где $f$ --- номер марки бумаги и $b_{f}$ --- время установки первой марки при $f = 1$ или время перехода с марки $f - 1$ на марку $f$; $U$ --- количество марок.
Порядок марок соответствует заданному.
Данный параметр отражает технологическое требование: смена марки бумаги на БДМ выполняется в фиксированной последовательности, определяемой свойствами сырья и настройками оборудования.

\[
O = \left\{ o_{1},\ \ldots,\ o_{Q} \right\}
\]
--- множество заказов, где каждый заказ описывается кортежем:
\[
o_{i} = \left( w_{i},\ q_{i},\ h_{i},\ f_{i},\ p_{i},\ l_{i},\ g_{i} \right),
\]
имеющим ширину $w_{i}$, количество $q_{i}$, срок $h_{i}$ (дата и время), номер марки бумаги $f_{i} \in [1,U]$, плотность $p_{i}$, количество слоёв $l_{i}$ и индикатор БРС $g_{i} \in \{ 0,1\}$; $Q$ --- количество заказов.
Каждый заказ полностью характеризуется набором параметров, необходимых для определения того, на каком оборудовании и из какого полуфабриката он может быть выполнен: ширина определяет схему раскроя, марка и плотность --- выбор тамбура, количество слоёв --- число тамбуров, одновременно подаваемых на ПРС, а индикатор $g_i$ --- необходимость использования бобино-резательного станка для узких форматов.
Параметр $l_i$ (число слоёв) отражает практику производства многослойной продукции (картон, гофра), когда на~ПРС одновременно подаётся несколько тамбуров; этот параметр непосредственно влияет на~число требуемых полуфабрикатов и~загрузку склада.
Индикатор $g_i$ разделяет заказы на~две группы: обычные ($g_i=0$), нарезаемые непосредственно из~тамбуров на~ПРС, и~малые ($g_i=1$), для~которых сначала на~ПРС формируются промежуточные рулоны, а~затем они нарезаются на~БРС.
Такое двухэтапное производство малых форматов характерно для~ЦБП \cite{sgp_ccp2021_cutting_stock_modes} и~в~явном виде ранее не~формализовалось в~моделях раскроя.

$H$ -- дата и время начала запуска производства.
Этот параметр фиксирует левую границу временного горизонта планирования; все операции расписания~$\overline{\overline{Z}}$ выполняются начиная с~момента~$H$.

$V$ -- размер кромки тамбура.
Кромка представляет собой технологический отход, срезаемый по~краям тамбура при~нарезке \cite{hinxman1980trimloss}; она уменьшает полезную ширину тамбура на~величину~$V$ и~учитывается в~ограничении допустимости раскроя.

Каждый тип оборудования описывается набором параметров, определяющих его производительность и технологические ограничения.

\[
M = \{ m_{1},\ \ldots,\ m_{I}\},\quad m_{i} = ({w^m}_{i},\ {v^m}_{i},\ {g^m}_{i})
\]
--- множество бумагоделательных машин (БДМ), где каждая БДМ имеет ширину ${w^m}_{i}$, время производства одного рулона ${v^m}_{i}$ и время смены плотности на 1~грамм ${g^m}_{i}$; $I$ --- количество БДМ.
Ширина БДМ определяет ширину выпускаемых тамбуров и тем самым задаёт верхнюю границу для схем раскроя на связанных станках.

\[
C = \{ c_{1},\ \ldots,\ c_{J}\},\quad c_{i} = ({w^c}_{i},\ {n^c}_{i},\ {v^c}_{i},\ {k^c}_{i},\ {g^c}_{i})
\]
--- множество продольно-резательных станков (ПРС), где каждый ПРС имеет максимальную ширину тамбура ${w^c}_{i}$, количество ножей ${n^c}_{i}$, время нарезки одного слоя для одного рулона ${v^c}_{i}$, время перестановки одного ножа ${k^c}_{i}$ и индикатор БРС ${g^c}_{i} \in \{ 0,1\}$; $J$ --- количество ПРС.
Параметр ${n^c}_{i}$ ограничивает максимальное число одновременно получаемых форматов в одном раскрое, а ${k^c}_{i}$ определяет штраф по времени за изменение конфигурации ножей между последовательными раскроями \cite{mobasher2013,martina2022, brucker2007_scheduling_algorithms}.
Индикатор ${g^c}_i$ объединяет ПРС и~БРС в~одном множестве~$C$, что позволяет единообразно описывать оба типа станков и~использовать одни и~те~же численные методы формирования планов нарезки.
Такой подход отличается от~моделей, рассматривающих только один тип резательного оборудования \cite{gilmore1961linear,haessler1991}.

\[
T = \{ t_{1},\ldots,\ t_{K}\},\quad t_{i} = ({r^T}_{i},\,{u^T}_{i})
\]
--- множество складов, где каждый склад имеет максимальную ширину тамбура ${r^T}_{i}$ и максимальное количество рулонов ${u^T}_{i}$; $K$ --- количество складов.
Склады играют роль буфера между этапами выпуска и нарезки: ограничение по ширине ${r^T}_{i}$ исключает размещение слишком широких тамбуров, а ограничение по количеству ${u^T}_{i}$ предотвращает переполнение и обеспечивает логистическую реализуемость плана.


\[
S = \{ s_{1},\ \ldots,\ s_{P}\}
\]
--- множество производственных помещений (цехов), где $P$ --- количество помещений.
Принадлежность оборудования помещениям задаётся входными параметрами:
${\rho^m}_{i} \in [1,P]$ --- номер помещения, в котором расположена БДМ $m_i$;
${\rho^c}_{j} \in [1,P]$ --- номер помещения ПРС/БРС $c_j$;
${\rho^T}_{k} \in [1,P]$ --- номер помещения склада $t_k$.
Данные параметры отражают компоновку производственных линий ЦБК: передача тамбуров и рулонов между стадиями производства, хранения и нарезки возможна только в пределах одного помещения, поэтому взаимодействующие БДМ, склады и ПРС/БРС должны принадлежать одному цеху (цеховое ограничение~10).
При $P=1$ ограничение вырождается и модель совпадает с постановкой без учёта помещений.

Включение складов в~модель является существенным расширением по~сравнению с~классическими постановками задачи раскроя, где полуфабрикат предполагается доступным в~неограниченном количестве \cite{gilmore1961linear,delorme2016bin}.
В~интегрированных постановках \cite{nonas2000combined,melega2018review,leao2016decomposition, pochet2006_prodplan_mip} склады моделируются через переменные запасов, однако в~настоящей работе состояние складов отслеживается на~уровне отдельных тамбуров ($\overline{T}$ и~$\overline{\overline{T}}$), что обеспечивает более точный контроль совместимости по~ширине, марке и~плотности.

Обозначим переменные решения, описывающие три взаимосвязанных компонента глобального плана: план производства тамбуров на БДМ, план нарезки на ПРС и расписание, связывающее их во времени.
Трёхкомпонентная структура решения $G=(\overline{\overline{X}},\overline{\overline{Y}},\overline{\overline{Z}})$ отражает реальную организацию производства, в~которой решения о~выпуске полуфабриката, формировании схем раскроя и~упорядочивании операций во~времени принимаются взаимосвязанно.
В~классических постановках \cite{gilmore1961linear,haessler1991} переменные решения ограничиваются планом нарезки~$\overline{\overline{Y}}$ (или объёмами использования схем раскроя); план производства~$\overline{\overline{X}}$ и~расписание~$\overline{\overline{Z}}$ либо не~рассматриваются, либо строятся независимо.
В~работах по~интеграции раскроя и~планирования партий \cite{nonas2000combined,poltroniere2016tema} добавляются переменные объёма выпуска, но~детальная структура расписания с~учётом состояния складов, как правило, не~моделируется.
Предлагаемый подход объединяет все три компонента, что позволяет проверять допустимость плана не~только по~ширинам и~объёмам, но~и~по~временным и~складским ограничениям.
Трёхкомпонентная структура решения важна и с прикладной точки зрения. Компонента $\overline{\overline{Y}}$ определяет, какие форматы и в каком сочетании должны быть получены на участке резки; из неё следует потребность в тамбурах определённых ширин, марок и плотностей, что задаёт содержание $\overline{\overline{X}}$. В свою очередь, уже построенные планы выпуска и нарезки нельзя считать реализуемыми, пока не проверена возможность их совместного выполнения во времени и по состоянию складов, что отражается в расписании $\overline{\overline{Z}}$. Следовательно, каждая из трёх компонент имеет самостоятельный смысл, но допустимое решение возникает только при их согласовании.

\[
\overline{\overline{X}} = \{{\overline{X}}_{1},\ldots,{\overline{X}}_{I}\}
\]
--- общий план производства бумаги.
${\overline{X}}_{i}$ --- план производства для $i$-ой БДМ, $i \in [1, I]$, 
состоящий из тамбуров $X_{i,j} \in {\overline{X}}_{i}$, 
которые необходимо произвести.
Каждый тамбур описывается кортежем:
\[
X_{i,j} = ({{w^X}}_{i,j},\ {{f^X}}_{i,j},\ {{p^X}}_{i,j}),
\]
имеющим ширину, совпадающую с шириной БДМ: ${{w^X}}_{i,j} = {w^m}_{i}$, номер марки бумаги ${f^X}_{i,j} \in [1,U]$ и плотность ${p^X}_{i,j}$.
План производства задаёт упорядоченную последовательность тамбуров, которые должны быть изготовлены на~каждой БДМ.
Равенство ${{w^X}}_{i,j}={w^m}_{i}$ отражает технологическое свойство бумагоделательной машины: ширина выпускаемого полуфабриката определяется конструкцией БДМ и~не~может быть произвольно изменена между выпусками.
Марка~${f^X}_{i,j}$ и~плотность~${p^X}_{i,j}$ назначаются в~соответствии с~текущими потребностями заказов и~подчиняются ограничению порядка марок (ограничение~1).
Тем самым компонента~$\overline{\overline{X}}$ определяет материальное обеспечение последующих операций нарезки: каждый тамбур из~$\overline{\overline{X}}$ после производства поступает на~склад и~становится доступным для~использования на~ПРС или~БРС.

\[
\overline{\overline{Y}} = \{{\overline{Y}}_{1},\ldots,{\overline{Y}}_{J}\}
\]
--- общий план нарезки бумаги.
${\overline{Y}}_{i}$ --- план нарезки для $i$-ой ПРС, $i \in [1,J]$, состоящий из раскроев $Y_{i,j} \in {\overline{Y}}_{i}$, которые необходимо выпустить.
Каждый раскрой описывается кортежем:
\[
Y_{i,j} = ({{y^Y}}_{i,j},\ {{f^Y}}_{i,j},\ {{p^Y}}_{i,j},\ {{l^Y}}_{i,j},\ {{g^Y}}_{i,j}),
\]
имеющим упорядоченный набор форматов $y_{i,j,k} \in {{y^Y}}_{i,j}$, где каждый формат:
\[
y_{i,j,k} = ({r^Y}_{i,j,k},\,{{g^Y}}_{i,j,k})
\]
состоит из ширины ${r^Y}_{i,j,k}$ и номера заказа ${{g^Y}}_{i,j,k} \in \{0\} \cup [1,Q]$ (если ${{g^Y}}_{i,j,k} = 0$, то нарезаемый формат~--- тамбур для БРС, а не конечный заказ);
номер марки бумаги ${{f^Y}}_{i,j} \in [1,U]$, плотность ${{p^Y}}_{i,j}$ и количество слоёв ${{l^Y}}_{i,j}$.
Каждый раскрой~$Y_{i,j}$ соответствует одному проходу тамбура (или нескольких тамбуров при~${{l^Y}}_{i,j}>1$) через станок~$c_i$: набор форматов~${{y^Y}}_{i,j}$ определяет, какие ширины одновременно получают из~полуфабриката за~один проход.
Разность между шириной использованного тамбура (с~учётом кромки~$V$) и~суммой ширин форматов в~${{y^Y}}_{i,j}$ составляет отход данного раскроя.
Нулевое значение номера заказа (${{g^Y}}_{i,j,k}=0$) означает, что получаемый рулон не~является конечной продукцией, а~предназначен для~дальнейшей нарезки на~БРС (станке с~${g^c}_j=1$); такая двухэтапная обработка характерна для~узких форматов.
Параметр~${{l^Y}}_{i,j}$ фиксирует число слоёв раскроя: при~${{l^Y}}_{i,j}>1$ на~ПРС одновременно подаётся несколько тамбуров одинаковой марки и~плотности, что требует их~одновременного наличия на~складе.

Количество перестановок ножей для $i$-ой ПРС и $j$-го плана различается при первом раскрое $j = 1$ и последующих раскроев $j > 1$.
При первом раскрое $j = 1$ необходимо выставить ножи ПРС перед первым форматом и после каждого формата:
\[
\sigma(i,j) = \left| y_{i,j} \right| + 1, \quad j = 1.
\]
Для последующих раскроев $j > 1$ количество ножей соответствует сумме предыдущих перестановок, разности количества форматов и количества совпадающих форматов в началах текущего и предыдущего раскроя:
\[
\sigma(i,j) = \sigma(i,j - 1) + \left| y_{i,j} \right| - \max_{n}\left\{ 0 \leq n \leq \min\left\{ \left| y_{i,j - 1} \right|,\left| y_{i,j} \right| \right\}:\ \forall k \leq n\ \ {r^Y}_{i,j - 1,k} = {r^Y}_{i,j,k} \right\},\quad j > 1.
\]
Физический смысл формулы состоит в~следующем.
На~ПРС ножи устанавливаются в~определённых позициях, задающих границы форматов.
При переходе от~раскроя~$j{-}1$ к~раскрою~$j$ ножи, уже~стоящие в~нужных позициях (совпадающий начальный префикс последовательности форматов), не~перемещаются; перестанавливаются лишь остальные.
Этот механизм стимулирует группировку схожих раскроев подряд: чем~длиннее общий префикс двух соседних раскроев, тем~меньше перестановок ножей и, следовательно, меньше время простоя оборудования.
Минимизация суммарного числа перестановок~$\tau(\overline{\overline{Y}})$ в~целевой функции непосредственно учитывает этот эффект.

\[
\overline{\overline{Z}}  = \{\overline{Z}_1,\,\overline{Z}_2,\,\ldots\}
\]
--- расписание, состоящее из операций $\overline{Z} \in \overline{\overline{Z}}$.
Расписание переводит структурные планы~$\overline{\overline{X}}$ и~$\overline{\overline{Y}}$ в~упорядоченную во~времени последовательность действий на~конкретных единицах оборудования.
Каждой операции соответствует либо выпуск одного тамбура на~БДМ, либо выполнение одного раскроя на~ПРС/БРС.
Помимо временных меток, операция хранит содержимое складов и~зарезервированное содержимое складов, что позволяет контролировать материальную обеспеченность нарезки и~предотвращать переполнение складов на~каждом шаге исполнения плана.
Каждая операция описывается кортежем:
\[
\overline{Z} = (h,\ \overline{h},\ e,\,d,\,z,\,\overline{T},\,\overline{\overline{T}},\ \overline{t'}),
\]
где:

\begin{itemize}
    \item $e\  = \ \{ 0,\ 1\}$ -- тип операции, при $e = 0$ -- это операция по выпуску тамбура на~БДМ, при $e = 1$ -- это операция по нарезке на~ПРС или~БРС.
    Разделение на~два типа операций позволяет единообразно описывать производственный цикл, в~котором выпуск полуфабриката и~его~нарезка чередуются и~выполняются на~разных единицах оборудования.
    \item $d$ -- номер машины, на~которой выполняется операция: при~$e=0$ это номер БДМ ($d\in[1,I]$), при~$e=1$~--- номер ПРС или~БРС ($d\in[1,J]$).
    $$d \in \left\{ 
    \begin{array}{r}
    I,e = 0 \\
    J,\ e = 1
    \end{array} 
    \right.\ $$
    \item $z$ -- порядковый номер элемента в~плане выпуска~$\overline{X}_d$ (при~$e=0$) или в~плане нарезки~$\overline{Y}_d$ (при~$e=1$) для~$d$-ой машины.
    Тем~самым тройка~$(e,d,z)$ однозначно идентифицирует конкретный тамбур или~раскрой, к~которому относится операция.
    
    \item $h$ -- время начала операции; определяется моментом готовности оборудования и~наличием необходимых ресурсов (при~$e=1$~--- тамбуров на~складе).
    \item $\overline{h}$ -- время окончания операции; зависит от~типа операции, характеристик оборудования и~параметров обрабатываемого полуфабриката.
При выпуске тамбура $(e=0)$ время конца $\overline{h}$ операции зависит от начала операции $h$, времени работы на БДМ ${v^m}_{d}$, плотности бумаги ${g^m}_{d}$ и скорости смены марки $b_f$ (если требуется замена), где $f={p^X}_{d,z}$ - номер марки.
Время окончания выпуска тамбура составляет:
    $$
    {
        \overline{h}=\left\{        
        \begin{array}{r}
                 h+{v^m}_{d}+{g^m}_{d}({p^X}_{d,z}-{p^X}_{d,z-1})+b_fI(f={f^X}_{d,z-1}),\\z>0,e=0\\
                 h+{v^m}_{d}+{g^m}_{d}{p^X}_{d,z}+b_f,z=0,e=0
        \end{array}
        \right.
    }
    $$.
При нарезке $(e=1)$ время зависит от скорости работы ПРС ${v^c}_{d}$, количества слоёв ${{l^Y}}_{d,z}$ и скорости ${k^c}_{d}$, количества $\sigma(d,z)$ перестановок ножей.
Время окончания нарезки рулона составляет:
    \begin{equation}
        \label{eq:slitting_time}
    {
        \overline{h}=\left\{        
        \begin{array}{r}
                 h+{v^c}_{z}{{l^Y}}_{d,z}+{k^c}_{d}(\sigma(d,z)-\sigma(d,z-1)),z>0,e=1\\
                 h+{v^c}_{z}{{l^Y}}_{d,z}+{k^c}_{d}\sigma(d,z),z=0,e=1
        \end{array}
        \right.
    }
    \end{equation}.
    
    \item $\overline{T} = \{{\overline{t}}_{1},\ldots,{\overline{t}}_{K}\}$ -- содержимое складов тамбуров, фиксирующее набор полуфабрикатов, физически находящихся на~складах в~момент выполнения данной операции.
$K$ -- количество складов.
Содержимое $i$-го склада состоит из набора тамбуров, полученных при выполнении планов по производству бумаги ${\overline{t}}_{i} = \{{\overline{t}}_{i,j}\}$.
${\overline{t}}_{i,j} = ({\overline{w^T}}_{i,j},\ {\overline{f^T}}_{i,j},\ {\overline{p^T}}_{i,j})$ -- тамбур на складе, имеющем ширину ${\overline{w^T}}_{i,j}$, марку ${\overline{f^T}}_{i,j}$ и плотность ${\overline{p^T}}_{i,j}$.
При операции по производству бумаги $(e = 0)$ рулон из плана не будет сразу включен в содержимое склада для текущей операции, так как будет произведен в будущем.
Когда наступит операция со временем начала $h$ не меньше времени окончания текущей операции $\overline{h}$, то данный рулон будет включен в склад $\overline{T}$.
При нарезке $(e = 1)$, тамбуры отправляемые на нарезку не включены.
Содержимое складов переходит к следующей операции, но без рулонов, отправленных на нарезку, и с новопроизведенными рулонами.
Таким образом, $\overline{T}$ отражает реально доступный набор полуфабрикатов в~момент начала операции и~используется для~проверки материальной обеспеченности нарезки (ограничение~7).
    \item $\overline{\overline{T}} = \{{\overline{\overline{t}}}_{1},\ldots,{\overline{\overline{t}}}_{K}\}$ -- зарезервированное содержимое складов тамбуров ${\overline{\overline{t}}}_{i} = \{{\overline{\overline{t}}}_{i,j}\}$, включающее содержимое $\overline{T} \subset \overline{\overline{T}}$ и~тамбуры, которые в~момент операции ещё находятся на~производстве БДМ, но~уже запланированы к~поступлению на~склад.
$K$ -- количество складов.
${\overline{\overline{t}}}_{i,j} = ({\overline{\overline{w^T}}}_{i,j},\ {\overline{\overline{f^T}}}_{i,j},\ {\overline{\overline{p^T}}}_{i,j})$ -- тамбур на складе, имеющем ширину ${\overline{\overline{w^T}}}_{i,j}$, марку ${\overline{\overline{f^T}}}_{i,j}$ и плотность ${\overline{\overline{p^T}}}_{i,j}$.
При операции по производству бумаги $(e = 0)$ рулон из плана будет включен в содержимое склада для текущей операции.
При нарезке $(e = 1)$, тамбуры отправляемые на нарезку не включены.
Зарезервированное содержимое складов переходит к следующей операции, но без рулонов, отправленных на нарезку.
Разделение складских переменных на~содержимое~$\overline{T}$ и~зарезервированное содержимое~$\overline{\overline{T}}$ отражает практическую ситуацию: в~момент планирования нарезки часть тамбуров уже находится на~складе и~готова к~использованию, а~часть~--- ещё производится на~БДМ и~поступит позже.
Ограничения по~складу (ограничения~8 и~9) применяются к~обеим переменным, что~позволяет предотвращать переполнение не~только в~текущий момент, но~и~с~учётом уже запущенных в~производство тамбуров.
Для двухэтапной нарезки это правило распространяется и на промежуточные рулоны с нулевым номером заказа: перед запуском производящей их операции резервируется место под выход после изъятия входных тамбуров. При завершении производства или резки ожидаемый полуфабрикат становится фактическим содержимым склада, оставаясь в общем зарезервированном содержимом; ёмкость освобождается только при изъятии полуфабриката на следующую обработку. Конечные заказные рулоны в этих складах полуфабриката не учитываются.
    
    \item $\overline{t'} = \{{\overline{t''}}_{i}\}$ -- множество тамбуров, изъятых со~склада для~выполнения операции нарезки ($e=1$); при производстве бумаги ($e=0$) используемые тамбуры отсутствуют: $\overline{t'} = \varnothing$.
${\overline{t''}}_{i} = ({\overline{w^t}}_{i},\ {\overline{f^t}}_{i},\ {\overline{p^t}}_{i})$ -- используемый тамбур, имеющий ширину ${\overline{w^t}}_{i}$, марку ${\overline{f^t}}_{i}$ и плотность ${\overline{p^t}}_{i}$.
\end{itemize}

\[
\overline{\overline{Y}} \;\xrightarrow{\text{определяет потребность}}\; \overline{\overline{X}} \;\xrightarrow[\text{согласование}]{\text{склад }\overline{T}}\; \overline{\overline{Z}}
\]

\[
G = (\overline{\overline{X}},\;\overline{\overline{Y}},\;\overline{\overline{Z}})
\]
--- глобальный план, состоящий из плана производства бумаги, плана нарезки бумаги и расписания.
Глобальный план является центральным объектом модели: он объединяет решения по выпуску полуфабриката, по формированию схем раскроя и по упорядочиванию операций во времени.
Задача состоит в нахождении такого плана $G$, который удовлетворяет всем ограничениям и минимизирует целевую функцию, определённую ниже.
Для более наглядной интерпретации этой структуры рассмотрим малый пример. Пусть имеется один заказ шириной $40$, один заказ шириной $20$, одна БДМ шириной $60$, один ПРС шириной $60$ и один склад. Тогда естественный план нарезки $\overline{\overline{Y}}$ может включать раскрой вида $(40,20)$, позволяющий из одного тамбура получить оба требуемых формата без дополнительного остатка сверх кромки. Такой раскрой определяет не только состав выпускаемой продукции, но и потребность в полуфабрикате: чтобы реализовать данную схему, необходимо иметь тамбур ширины $60$, совместимый с параметрами выбранной БДМ.
Соответственно, в компоненте $\overline{\overline{X}}$ появляется выпуск одного тамбура ширины $60$ требуемой марки и плотности. Однако даже после задания $\overline{\overline{X}}$ и $\overline{\overline{Y}}$ решение ещё не является полностью определённым. Необходимо зафиксировать, когда именно БДМ изготовит этот тамбур, когда он поступит на склад и в какой момент будет подан на ПРС для выполнения раскроя. Эта информация задаётся расписанием $\overline{\overline{Z}}$, которое превращает структурное описание плана в выполнимый производственный сценарий.
Из примера видно, что без компоненты $\overline{\overline{Z}}$ даже простой набор решений нельзя считать реализуемым в строгом смысле. Формально можно указать раскрой и выпуск тамбура, но без временного согласования остаются неучтёнными доступность полуфабриката, занятость оборудования и прохождение материала через склад. Именно поэтому в настоящей работе глобальный план трактуется как единый объект, а не как механическая сумма нескольких независимых подзадач.
Прикладная ценность такого представления состоит в том, что каждая компонента глобального плана отвечает на отдельный производственный вопрос, но только совместно они образуют решение, пригодное для передачи в работу. Компонента $\overline{\overline{Y}}$ отвечает за структуру раскроя и использование ширины полуфабриката, компонента $\overline{\overline{X}}$ --- за обеспечение этой структуры материальным ресурсом, а компонента $\overline{\overline{Z}}$ --- за календарную реализуемость на доступном наборе оборудования. Если хотя бы одна из этих частей отсутствует или формируется независимо от остальных, предприятие получает не план производства, а лишь частичное описание отдельных его аспектов.
Именно поэтому математическая модель далее строится не вокруг одной «главной» подзадачи, а вокруг механизма согласования трёх компонент. В терминах исследования это означает, что допустимость решения определяется не локальным выполнением какого-либо одного ограничения, а совместимостью всех структурных, материальных и временных характеристик плана. Такая интерпретация особенно важна для последующего раздела о численных методах, где любые изменения в раскроях должны автоматически отражаться на выпуске полуфабриката и расписании.

\section{Целевая функция}\label{sec:model/objective}
Оптимальность плана можно измерить по следующим критериям.
Порядок перечисления соответствует приоритету критериев:
\begin{enumerate}
    \item суммарная ширина отходов;
    \item максимальное количество раскроев из планов для каждой ПРС;
    \item общее количество перестановок ножей.
\end{enumerate}

Далее вводится векторная целевая функция $\mu(\overline{\overline{Y}})$ и даётся интерпретация её компонент, соответствующих указанным критериям.


% --- начало расширенного описания целевой функции ---
Векторная целевая функция
\[
\mu(\overline{\overline{Y}}) = \Bigl\langle \omega(\overline{\overline{Y}}),\, \varphi(\overline{\overline{Y}}),\, \tau\bigl(\overline{\overline{Y}}\bigr) \Bigr\rangle \rightarrow \min
\]
отражает три ключевых аспекта качества плана нарезки $\overline{\overline{Y}}$.

Суммарная ширина отходов определяется как разность между шириной тамбура и суммарной шириной форматов в каждом раскрое каждой ПРС:
\[
\omega(\overline{\overline{Y}})
\;=\;
\sum_{i=1}^{J}\sum_{j=1}^{|\overline{Y}_i|}
\left(
\,l^Y_{i,j}\,{w^c}_{i}
\;-\;
\sum_{k=1}^{|y^Y_{i,j}|} r^Y_{i,j,k}
\right),
\]
где ${w^c}_{i}$ --- ширина тамбура, используемого на $i$-й ПРС, $l^Y_{i,j}$ --- число слоёв в $j$-м раскрое, а $r^Y_{i,j,k}$ --- ширина $k$-го формата в этом раскрое.
Этот критерий оценивает суммарный отход материала и включает в себя кромку; его минимизация обеспечивает экономию сырья и снижение экологических издержек \cite{haessler1991,mobasher2013}.
В многокритериальных работах по задаче раскроя минимизация потерь часто рассматривается как основной показатель качества решения \cite{araujo2014genetic}.

Вторая компонента
\[
\varphi(\overline{\overline{Y}})\;=\;\max_{j=1,\ldots,J} |\overline{Y}_j|
\]
определяет максимальное число раскроев на одной ПРС.
Число раскроев для каждой машины прямо пропорционально трудоёмкости и времени выполнения заказов.
Ограничение пиковых нагрузок позволяет равномернее распределить работу между машинами, снизить время ожидания и обеспечить более стабильную эксплуатацию оборудования \cite{cerqueira2021,vanderbeck2000}.

Третья компонента
\[
\tau\bigl(\overline{\overline{Y}}\bigr)\;=\;\sum_{i=1}^{J} \sigma\bigl(i,\,|\overline{Y}_i|\bigr),
\]
где $\sigma(i,j)$ — количество перестановок ножей при $j$‑м раскрое на $i$‑й ПРС, отражает общую трудоёмкость переналадок.
Перестановка ножей требует времени и сопровождается простоем, поэтому минимизация $\tau(\overline{\overline{Y}})$ помогает уменьшить длительность производственного цикла \cite{cyberleninka_slitting_params,mobasher2013}.

Вычислительная сложность однократного вычисления целевой функции $\mu(\overline{\overline{Y}})$ определяется стоимостью вычисления каждой из~трёх компонент.
Обозначим $N_Y=\sum_{j=1}^{J}|\overline{Y}_j|$~--- суммарное число раскроев во~всём плане нарезки, $R=\sum_{i=1}^{Q}q_i$~--- суммарное число рулонов, $L_{\max}=\max_{j,k}l^Y_{j,k}$~--- максимальное число слоёв в~одном раскрое, $n^c_{\max}=\max_j n^c_j$~--- максимальное число ножей.
Вычисление $\omega(\overline{\overline{Y}})$ требует обхода всех $N_Y$ раскроев плана нарезки~$\overline{\overline{Y}}$: для каждого раскроя суммируются ширины полуфабрикатов (до~$L_{\max}$ слоёв) и~форматов; суммарно $O(N_Y \cdot L_{\max} + R)$ операций.
Вычисление $\varphi(\overline{\overline{Y}})$ сводится к~нахождению максимума по~$J$ станкам: $O(J)$.
Вычисление $\tau(\overline{\overline{Y}})$ требует последовательного расчёта $\sigma(i,j)$ для~каждого раскроя; сравнение двух смежных раскроев выполняется за~$O(n^c_{\max})$, суммарно $O(N_Y \cdot n^c_{\max})$.
Таким образом, общая сложность вычисления целевой функции составляет:
\begin{equation}\label{eq:objective_complexity}
T_{\mu} = O\bigl(N_Y \cdot (L_{\max} + n^c_{\max}) + R\bigr).
\end{equation}
Эта оценка используется при~анализе вычислительной сложности численных методов в~главе~\ref{ch:algorithm}.

Компоненты $\mu(\overline{\overline{Y}})$ могут быть интерпретированы в терминах производственных показателей.
Величина $\omega(\overline{\overline{Y}})$ выражается в единицах ширины (в используемых единицах измерения входных данных) и на практике может быть дополнительно приведена к относительному показателю, например к доле отхода в процентах от суммарной ширины использованного полуфабриката.
Показатель $\varphi(\overline{\overline{Y}})$ измеряется количеством операций нарезки (раскроев) и отражает загрузку участка резки; при прочих равных уменьшение $\varphi(\overline{\overline{Y}})$ снижает риск возникновения «пиков» по времени выполнения и уменьшает число технологических переходов.
Показатель $\tau(\overline{\overline{Y}})$ измеряется количеством перестановок ножей и напрямую связан с простоем оборудования и трудоёмкостью переналадки.

В прикладной интерпретации данные показатели часто рассматриваются совместно: план с минимальным отходом может приводить к большому числу раскроев и переналадок (что ухудшает производительность), тогда как план с упрощённой последовательностью схем уменьшает переналадки, но может увеличивать отход.
Поэтому сравнение альтернативных планов удобно выполнять как по формальным критериям, так и по производственным показателям (например, доля отхода, количество операций на смену, суммарное время переналадок).

В теории многокритериальной оптимизации известно несколько подходов к работе с векторной целевой функцией \cite{васильев2016методы,зенков2017методы,струченков2017методы, deb2001_moo_ea, tkindt2006_multicriteria_scheduling}.
Применительно к~$\mu(\overline{\overline{Y}})$ можно выделить следующие основные методы.
\begin{itemize}
  \item \emph{Лексикографическая минимизация.}
Критерии упорядочиваются по важности, и решение выбирается по принципу последовательного улучшения: сначала минимизируется $\omega(\overline{\overline{Y}})$, затем при равных значениях первого критерия — $\varphi(\overline{\overline{Y}})$, и только в~случае равенства первых двух — $\tau(\overline{\overline{Y}})$.
Такой подход оправдан, когда приоритет одного критерия безусловно доминирует: в задачах раскроя уменьшение отходов обычно имеет наибольшее экономическое значение \cite{haessler1991,cerqueira2021,vanderbeck2000}.

  \item \emph{Взвешенная сумма (линейная свёртка).}
Критерии комбинируются в скалярную функцию
$\mu_{\text{ws}}(\overline{\overline{Y}}) = \omega(\overline{\overline{Y}}) + \alpha\,\varphi(\overline{\overline{Y}}) + \beta\,\tau(\overline{\overline{Y}})$
с~весовыми коэффициентами $\alpha,\beta>0$.
Метод применяется в генетических алгоритмах для ЗЛР, где веса задают компромисс между расходом материала и производственными издержками \cite{araujo2014genetic,golfeto2009genetic};
в~интегрированных постановках «раскрой + планирование партий» аналогичная линейная комбинация объединяет затраты на производство, хранение и потери при раскрое \cite{poltroniere2016tema,nonas2000combined}.
Основной недостаток — чувствительность результата к~выбору весов, которые обычно определяются экспертом или по итогам вычислительных экспериментов. Для рассматриваемой задачи это особенно существенно, поскольку пришлось бы заранее задавать относительную «цену» единицы отхода, одного дополнительного раскроя и одной перестановки ножей. На практике такая калибровка часто оказывается нестабильной: значения критериев имеют разную размерность и по-разному проявляются на различных наборах заказов, поэтому единый набор весов может хорошо работать в одном сценарии и приводить к нежелательным компромиссам в другом.

  \item \emph{Метод $\varepsilon$-ограничений.}
Один из критериев (как правило, главный) оптимизируется, а~остальные ограничиваются сверху заданными порогами:
$\omega(\overline{\overline{Y}}) \to \min$ при $\varphi(\overline{\overline{Y}}) \le \varepsilon_1$, $\tau(\overline{\overline{Y}}) \le \varepsilon_2$.
Метод позволяет строить множество Парето-оптимальных решений путём варьирования $\varepsilon_1,\varepsilon_2$ и не~требует назначения весов \cite{васильев2016методы}.
В~задачах раскроя подобный подход реализуется, когда число раскроев или перестановок ножей жёстко ограничивается технологическими нормативами \cite{mobasher2013,vanderbeck2000}. Для прикладного использования метода это означает необходимость заранее определить допустимые пороги для второстепенных критериев. В автоматизированной системе такое требование не всегда удобно, поскольку эти пороги зависят от текущей загрузки оборудования, состава заказов и организационных приоритетов предприятия.

  \item \emph{Целевое программирование.}
Для каждого критерия задаётся целевое значение $\omega^*$, $\varphi^*$, $\tau^*$, и минимизируется суммарное или максимальное отклонение от целей.
Метод удобен, когда предприятие формулирует требования в~форме плановых нормативов \cite{зенков2017методы}.

  \item \emph{Минимаксная (Чебышёвская) свёртка.}
Минимизируется наибольшее взвешенное относительное отклонение от~идеальной точки:
$\mu_{\infty}(\overline{\overline{Y}}) = \max\!\bigl\{\lambda_1\,\delta\omega(\overline{\overline{Y}}),\;\lambda_2\,\delta\varphi(\overline{\overline{Y}}),\;\lambda_3\,\delta\tau(\overline{\overline{Y}})\bigr\} \to \min$,
где $\delta(\cdot)$ — нормированное отклонение от наилучшего достижимого значения.
Такой подход гарантирует Парето-оптимальность найденного решения даже при невыпуклом фронте Парето \cite{васильев2016методы,струченков2017методы}.

  \item \emph{Парето-оптимальность и эволюционные методы.}
 Вместо скаляризации строится множество Парето-оптимальных (недоминируемых) решений.
 В~эволюционных алгоритмах отбор особей осуществляется на основе доминирования: решение $\overline{\overline{Y}}'$ доминирует $\overline{\overline{Y}}''$, если оно не хуже по всем критериям и строго лучше хотя бы по одному \cite{goldberg1989genetic,gladkov2010, deb2001_moo_ea, eiben2015_intro_ec}.
 Такие подходы применялись к~задачам раскроя с~несколькими целями — в~генетических \cite{golfeto2009genetic,sugianto2022rebar} и других метаэвристических алгоритмах \cite{lai1997developing}.
 Итоговое решение выбирается лицом, принимающим решение, из множества несравнимых альтернатив. Для программного комплекса, ориентированного на автоматический выпуск одного плана по входному набору данных, это создаёт дополнительную проблему: после завершения оптимизации требуется ещё одна процедура выбора окончательного варианта, которая выходит за рамки собственно численного метода.
\end{itemize}

Выбор конкретного способа агрегирования критериев зависит от требований предприятия и особенностей производства.
В~настоящей работе используется лексикографическое упорядочение критериев.
 Этот выбор обусловлен следующими соображениями:
(1)~в~ЦБП минимизация отходов материала является безусловным приоритетом с~наибольшим экономическим эффектом;
(2)~лексикографический порядок не~требует назначения весов и~обеспечивает детерминированное сравнение альтернативных планов;
(3)~критерии $\varphi(\overline{\overline{Y}})$ и $\tau(\overline{\overline{Y}})$ играют роль вторичных показателей, улучшаемых лишь при фиксированном уровне отхода, что соответствует практике предприятий отрасли \cite{haessler1991,melega2018review}.
 Дополнительно выбор лексикографического порядка обусловлен тем, что критерии имеют различную экономическую природу и измеряются в несопоставимых единицах. Значение $\omega(\overline{\overline{Y}})$ непосредственно отражает потери материала и потому интерпретируется как главный показатель эффективности. Показатель $\varphi(\overline{\overline{Y}})$ характеризует интенсивность загрузки участка резки и важен прежде всего для обеспечения организационной выполнимости плана. Показатель $\tau(\overline{\overline{Y}})$ фиксирует трудоёмкость переналадок и влияет на длительность простоев оборудования. Использование единой взвешенной суммы потребовало бы предварительно выбрать коэффициенты, отражающие относительную «цену» единицы отхода, одного дополнительного раскроя и одной перестановки ножа; такие коэффициенты зависят от конкретного предприятия и могут существенно меняться от периода к периоду. Лексикографический подход позволяет избежать этой неопределённости и одновременно сохранить прозрачную производственную интерпретацию выбора лучшего решения.
 Важно и то, что лексикографическое сравнение обеспечивает воспроизводимость результатов вычислительных экспериментов и программной реализации. При одинаковом наборе допустимых планов правило выбора не зависит от дополнительной экспертной настройки и даёт однозначный ответ, какой из двух планов считается лучшим. Это упрощает как реализацию операторов сравнения в программном комплексе, так и интерпретацию результатов при сравнении полного перебора и генетического поиска. По существу, именно лексикографический порядок в наибольшей степени соответствует задаче автоматизированного расчёта: он не требует заранее назначать веса, не требует задавать допустимые пороги и не перекладывает выбор окончательного плана на пользователя после завершения вычислений.
 
 % --- конец расширенного описания целевой функции ---
 
 \section{Ограничения}
 Введённые ограничения можно интерпретировать как четыре взаимосвязанные группы требований. Первая группа описывает технологические правила производства полуфабриката, прежде всего порядок марок бумаги на БДМ. Вторая группа относится к формированию допустимых раскроев и выполнению заказов по объёму. Третья группа связывает планы выпуска и нарезки с календарной выполнимостью, включая наличие операций в расписании и соблюдение сроков. Четвёртая группа описывает логистическую реализуемость решения за счёт контроля состояния складов и доступности полуфабриката в моменты нарезки. Такое представление важно, поскольку подчёркивает, что ограничения не являются независимыми: нарушение складского или временного условия может сделать недопустимым план, формально корректный по ширинам и объёмам.
 \begin{enumerate}
     \item Порядок марок в производстве тамбуров должен соответствовать заданному:
     \[
     \forall\, i,\,j:\quad {{f^X}}_{i,j} \leq {{f^X}}_{i,j + 1}.
    \]

    Данное ограничение фиксирует технологически заданную очередность марок бумаги и исключает решения, в которых последовательность марок нарушается.

    \item Для раскроя его сумма ширин заказов и кромки тамбура не должна превышать ширину ПРС:
    \[
    V + \sum_{k = 1}^{\left| y_{i,j} \right|}{r^Y}_{i,j,k} \leq {w^c}_{i},\quad i=\overline{1,Q},\; j=\overline{1,U}.
    \]
Общий план нарезки бумаги должен удовлетворять всем заказам:
    \[
    \forall\, n \in [1,Q]:\quad \sum_{i,j,k} I\bigl( {{g^Y}}_{i,j,k} = n \bigr) = q_{n}.
    \]

    Первая часть ограничивает допустимость каждого раскроя по ширине (с учётом кромки), а вторая гарантирует выполнение объёмов всех заказов в рамках общего плана нарезки.

    \item Расписание определяет последовательность выполнения планов производства бумаги и нарезки, поэтому количество операций равно сумме всех планов для каждой БДМ и ПРС:
    \[
    \left| \overline{\overline{Z}} \right| = \sum_{i = 1}^{I}\left| {\overline{X}}_{i} \right| + \sum_{i = 1}^{J}\left| {\overline{Y}}_{i} \right|.
    \]

    Это условие обеспечивает полноту расписания: каждой операции выпуска и каждой операции нарезки соответствует запись в расписании, и тем самым план допускает временную интерпретацию.

    \item Для каждого заказа срок должен быть не раньше, чем время конца операции, в которой есть нарезка соответствующего заказа:
    \[
    \forall\,\overline{Z},\;\forall\, i \in [1,Q]:\quad \overline{h} \leq h_{i}.
    \]

    Ограничение связывает план со сроками исполнения: выполнение заказов допускается только в том случае, если соответствующие операции нарезки завершаются не позже заданных сроков.

    \item Для любого тамбура из плана по производству бумаги должна быть соответствующая операция в расписании:
    \[
    \forall\, i \in [1, I],\; j \in [1,|{\overline{X}}_{i}|]:\quad \exists\,\overline{Z} \in \overline{\overline{Z}}\ \text{с} \ e = 0,\; d = i,\; z = j.
    \]

    Данное условие исключает ситуации, когда тамбур присутствует в плане выпуска, но не включён в расписание, то есть фактически не производится.

    \item Для любого раскроя из плана по нарезке бумаги должна быть соответствующая операция в расписании:
    \[
    \forall\, i \in [1, J],\; j \in [1,|{\overline{Y}}_{i}|]:\quad \exists\,\overline{Z} \in \overline{\overline{Z}}\ \text{с} \ e = 1,\; d = i,\; z = j.
    \]

    Аналогично предыдущему, это ограничение гарантирует, что каждая операция нарезки из плана действительно запланирована во времени.

    \item Для осуществления нарезки ($e = 1$) должны присутствовать необходимые тамбуры на складе в количестве, равном числу слоёв:
    \[
    \forall\,\overline{Z}:\quad \sum_{i,j} I\!\left( \sum_{k}{r^Y}_{d,z,k} + V \leq {\overline{w^T}}_{i,j} \;\&\; {{f^Y}}_{d,z} = {\overline{f^T}}_{i,j} \;\&\; {{p^Y}}_{d,z} = {\overline{p^T}}_{i,j} \right) = {{l^Y}}_{d,z}.
    \]

    Ограничение обеспечивает материальную выполнимость операции нарезки: к моменту операции на складе должны находиться подходящие тамбуры, совместимые по ширине, марке и плотности, причём их число соответствует требуемому числу слоёв.

    \item В содержимом складов тамбуров $\overline{T}$ тамбуры не должны превышать ограничения соответствующего склада по ширине и по количеству:
    \[
    {\overline{w^T}}_{i,j} \leq {r^T}_{i}\;\forall\,{\overline{t}}_{i,j} \in {\overline{t}}_{i},\qquad \left| {\overline{t}}_{i} \right| \leq {u^T}_{i}.
    \]

    Данное условие ограничивает содержимое склада: на складе не может находиться полуфабрикат шире допустимого и сверх заданной ёмкости.

    \item Для запланированного содержимого складов $\overline{\overline{T}}$ тамбуры не должны превышать ограничения соответствующего склада по ширине и по количеству:
    \[
    {\overline{\overline{w^T}}}_{i,j} \leq {r^T}_{i}\;\forall\,{\overline{\overline{t}}}_{i,j} \in {\overline{\overline{t}}}_{i},\qquad \left| {\overline{\overline{t}}}_{i} \right| \leq {u^T}_{i}.    
    \]
 
    Это ограничение распространяет требования вместимости на зарезервированное содержимое склада, включая полуфабрикаты, которые уже запущены в производство и будут поступать на склад в будущем.

    \item Взаимодействующее оборудование должно находиться в одном производственном помещении (цеховое ограничение).
    Тамбур, произведённый на БДМ $m_d$ (операция с~$e=0$), может поступать только на склад, расположенный в том же помещении:
    \[
    \forall\,\overline{Z}\ \text{с}\ e = 0:\quad {\overline{\overline{t}}}_{i,j} \ \text{--- новый тамбур операции} \;\Rightarrow\; {\rho^T}_{i} = {\rho^m}_{d}.
    \]
    При нарезке на ПРС/БРС $c_d$ (операция с~$e=1$) используемые тамбуры могут изыматься только со складов того же помещения:
    \[
    \forall\,\overline{Z}\ \text{с}\ e = 1,\ \forall\,{\overline{t''}} \in \overline{t'}:\quad {\overline{t''}} \in {\overline{t}}_{i} \;\Rightarrow\; {\rho^T}_{i} = {\rho^c}_{d}.
    \]
    Цеховое ограничение исключает планы, в которых полуфабрикат перемещается между оборудованием разных помещений: материальный поток «БДМ $\to$ склад $\to$ ПРС/БРС» замыкается внутри одного цеха. Совместно с ограничением~7 оно сужает множество тамбуров, пригодных для конкретной операции нарезки, и тем самым влияет на допустимость расписания и распределение заказов между машинами.
\end{enumerate}
Ограничения можно дополнительно сгруппировать по их роли в обеспечении выполнимости плана. Первые условия отвечают за технологическую допустимость: они задают порядок марок и корректность использования оборудования, тем самым исключая решения, противоречащие базовым правилам производственного процесса. Следующая группа условий обеспечивает временную и структурную полноту плана: она требует, чтобы каждому элементу выпуска и каждому раскрою соответствовала операция в расписании и чтобы выполнение заказов укладывалось в заданные сроки.
Отдельное место занимает ограничение материальной обеспеченности резки. Оно не позволяет назначить раскрой во времени, если к соответствующему моменту отсутствуют необходимые тамбуры требуемой ширины, марки, плотности и кратности слоёв. Именно это условие превращает склад из пассивного хранилища в активный элемент модели, влияющий на допустимость решения.
Наконец, две последние группы ограничений отвечают за логистическую реализуемость плана. Они контролируют не только содержимое складов, но и зарезервированное содержимое складов, предотвращая переполнение и размещение полуфабриката, несовместимого с параметрами конкретного склада. К этой же группе относится цеховое ограничение (п.~10): оно гарантирует, что материальный поток между БДМ, складами и ПРС/БРС не выходит за пределы одного производственного помещения. Тем самым каждая группа ограничений отсекает свой класс некорректных решений, а их совместное действие обеспечивает переход от абстрактной схемы раскроя к выполнимому глобальному плану.

Для пояснения смысла введённых структур рассмотрим небольшой пример.
Пусть $I=1$ (одна БДМ), $J=1$ (один станок резки), $K=1$ (один склад), а множество заказов $O$ содержит несколько элементов, различающихся по ширине $w_i$ и объёму $q_i$.
Тогда план выпуска $\overline{\overline{X}}$ задаёт последовательность тамбуров $X_{1,j}$, которые необходимо изготовить, план нарезки $\overline{\overline{Y}}$ задаёт набор раскроев $Y_{1,j}$ (каждый раскрой определяет, какие форматы $y_{1,j,k}$ одновременно получают из одного тамбура), а расписание $\overline{\overline{Z}}$ фиксирует порядок выполнения выпусков ($e=0$) и нарезок ($e=1$) с учётом доступности полуфабрикатов на складе $\overline{T}$.

В таком примере глобальный план $G = (\overline{\overline{X}}, \overline{\overline{Y}}, \overline{\overline{Z}})$ можно интерпретировать как цепочку решений: сначала формируется $\overline{\overline{X}}$ --- какие тамбуры и в какой последовательности изготовить, затем $\overline{\overline{Y}}$ --- какие схемы раскроя применить, и наконец $\overline{\overline{Z}}$ --- когда направлять тамбуры на резку и какие партии накопить на складе.
При этом ключевым является баланс между компонентами $\mu(\overline{\overline{Y}})$: более «плотные» раскрои уменьшают $\omega(\overline{\overline{Y}})$, но могут увеличивать $\varphi(\overline{\overline{Y}})$ и $\tau(\overline{\overline{Y}})$.

При практическом применении модели полезно проверять выполнимость полученного плана по набору инвариантов.
Такие проверки позволяют выявлять ошибки данных и некорректные решения на раннем этапе.
\begin{itemize}
    \item Выполнение заказов: объёмы выпуска по каждому заказу должны соответствовать требованиям (каждый заказ удовлетворён в полном объёме).
    \item Выполнимость раскроев: для каждой операции нарезки соблюдены ограничения по ширине и по технологическим параметрам станка.
    \item Наличие полуфабриката: операции нарезки выполняются только при наличии требуемых тамбуров на складе (с учётом числа слоёв).
    \item Непереполнение складов: содержимое и зарезервированное содержимое складов не нарушает ограничений ёмкости и допустимой ширины.
    \item Соблюдение цехового ограничения: каждая операция связывает только оборудование и склады, расположенные в одном производственном помещении.
    \item Соблюдение сроков: все операции, относящиеся к заказам, завершаются не позже заданных сроков.
\end{itemize}
Указанные проверки не изменяют модель, но обеспечивают интерпретируемость и воспроизводимость результатов, что важно при реализации и экспериментальной валидации. С практической точки зрения набор инвариантов выполняет роль содержательной спецификации корректного решения. Он показывает, что глобальный план должен быть не просто оптимальным по значению целевой функции, но и внутренне непротиворечивым: все заказы должны быть обеспечены, каждая операция должна иметь материальное основание, а складские и временные ограничения не должны нарушаться ни на одном шаге исполнения.
Такой подход важен ещё и потому, что в интегрированных постановках ошибка часто проявляется не в виде явного нарушения одной формулы, а как несогласованность между несколькими компонентами решения. Например, план может выглядеть корректным по отходам и числу раскроев, но оказаться невыполнимым из-за отсутствия полуфабриката в нужный момент или из-за кратковременного переполнения склада. Поэтому интерпретация ограничений через набор проверяемых инвариантов облегчает переход от формальной модели к программной реализации и делает результаты вычислительных экспериментов более прозрачными для анализа.

Таким образом, дальнейшее изложение сводится к задаче оптимизации: требуется построить глобальный план, обеспечивающий минимальное значение целевой функции при выполнении ограничений, то есть найти $G$ такое, что $\mu(\overline{\overline{Y}}) \rightarrow \min$.

Рассматриваемая постановка относится к классу трудных комбинаторных задач.
Уже классическая задача линейного раскроя (минимизация числа заготовок для заданного набора форматов) является NP-трудной \cite{garey1979computers}.
Добавление каждого нового аспекта~--- очерёдности марок, многослойных заказов, складских ограничений, переналадок и~временных рамок~--- не~уменьшает сложность, поскольку исходная задача является частным случаем расширенной при~$U=1$, $K=0$, $l_i=1$, $P=1$ и~отсутствии ограничений по~срокам.
Более того, необходимость согласования трёх компонент плана ($\overline{\overline{X}}$, $\overline{\overline{Y}}$, $\overline{\overline{Z}}$) порождает дополнительные зависимости: изменение плана нарезки влияет на~требуемый набор тамбуров и, как следствие, на~допустимость расписания с~учётом складов.
Эта взаимосвязь исключает возможность независимой оптимизации компонент и~обосновывает применение метаэвристических методов, описанных в~главе~\ref{ch:algorithm}.


С практической точки зрения размерность определяется количеством заказов $Q$, числом марок бумаги $U$, количеством оборудования $I$ (БДМ) и $J$ (станки резки), а также длиной планов $|\overline{X}_i|$ и $|\overline{Y}_i|$.
Для типичных промышленных экземпляров ($Q\sim 50$--$200$, $I\sim 2$--$5$, $J\sim 3$--$10$) полный перебор неприменим, что подтверждается экспоненциальной оценкой $O(J^R)$ (см.~раздел~\ref{sec:alg/full_enum}).

Построенная модель задаёт пространство исходно допустимых решений и целевые критерии для оптимизации плана нарезки. Далее рассматривается исполнение уже сформированных планов при случайных возмущениях. Исходная лексикографическая оценка не заменяет оценку устойчивости: нарушение сроков в имитационном прогоне регистрируется как результат воздействия возмущений, а не устраняется остановкой прогона в момент первого опоздания.

\section{Имитационное моделирование устойчивости исполнения плана}\label{sec:model/simulation}\label{sec:alg/simulation}
Входом дискретно-событийной модели является фиксированный глобальный план $G=(\overline{\overline{X}},\overline{\overline{Y}},\overline{\overline{Z}})$, удовлетворяющий исходным ограничениям. Имитация не строит новое расписание и не изменяет раскрои, назначения операций машинам или порядок операций на каждой машине. Меняются фактические моменты запуска и завершения, определяемые возмущениями, наличием материала и свободной складской ёмкости. Допускается ожидание ресурсов, но не обход предыдущей операции на той же машине и не перенос операции на другое оборудование.

Обозначим через $\mathcal{U}$ список операций исходного расписания, упорядоченный по плановому началу с фиксированным разрешением совпадений. Для $u\in\mathcal{U}$ известны тип $e(u)\in\{0,1\}$, номер машины $d(u)$, начало $h(u)$, конец $\overline{h}(u)$ и длительность $\pi(u)=\overline{h}(u)-h(u)>0$. Машину будем идентифицировать парой $a=(e,d)$, чтобы одинаковые номера БДМ и ПРС не смешивались. Модельное время обозначим $t^{*}$, фактическое первое начало и окончание операции~--- $h^{*}(u)$ и $\overline{h}^{*}(u)$.

Набор параметров производственной среды задаётся вектором
\[
\Theta=(\Delta^{x},\Delta^{y},\nu^{x},\nu^{y},\theta),
\qquad 0\leq\Delta^{x},\Delta^{y}<1,\quad \nu^{x},\nu^{y}\geq0,\quad\theta>0.
\]
Здесь $\Delta^{x},\Delta^{y}$~--- границы относительного разброса длительностей производства и резки, $\nu^{x},\nu^{y}$~--- интенсивности отказов соответствующего оборудования за единицу времени работы, $\theta$~--- среднее время восстановления. Отдельно задаются параметры эксперимента: число повторений $N_{\Theta}$ и начальное состояние генератора $\zeta$. Они управляют точностью и воспроизводимостью оценки, но не характеризуют производственную среду.

При первом запуске операции разыгрывается требуемое время работы без ремонтов:
\[
\widetilde{\pi}(u)=\pi(u)(1+\xi_u),\qquad
\xi_u\sim U(-\Delta^{(e(u))},\Delta^{(e(u))}),
\qquad \Delta^{(0)}=\Delta^x,\quad\Delta^{(1)}=\Delta^y.
\]
Наработка до отказа на каждом непрерывном участке работы и длительность ремонта моделируются как
\[
\chi\sim\mathrm{Exp}(\nu^{(e)}),\qquad
\widetilde{\theta}\sim\mathrm{Exp}(1/\theta),\qquad
\nu^{(0)}=\nu^x,\quad\nu^{(1)}=\nu^y.
\]
При $\nu^{(e)}=0$ отказ не разыгрывается, что эквивалентно $\chi=+\infty$. Длительности, отказы и ремонты в базовой вероятностной модели независимы. Отказы учитываются только во время выполнения операций, а не в периоды ожидания или ремонта. После восстановления продолжается оставшаяся часть той же операции без повторного расходования сырья и без потери уже выполненной работы; время до следующего отказа разыгрывается заново. Поэтому $\widetilde{\pi}(u)$ не включает ремонты, тогда как интервал $\overline{h}^{*}(u)-h^{*}(u)$ включает все её приостановки.

Это приближение описывает возобновляемую обработку и не моделирует брак, повторное изготовление после отказа и общие для нескольких машин причины остановок. Исходная длительность уже учитывает предусмотренные математической постановкой технологические переходы; отдельный случайный календарь переналадок в данной базовой модели не вводится.

Состояние имитационного прогона, материальные потоки и~события описываются следующим образом. Состояние прогона включает $t^{*}$, статусы операций (не запущена, выполняется, приостановлена, завершена), статусы машин, остатки времени обработки, фактические склады $\overline{T}$, их зарезервированное содержимое $\overline{\overline{T}}$ и список будущих событий $\mathcal{F}$. Список можно хранить в двоичной куче с минимальным ключом времени. События включают завершение, отказ и восстановление; дополнительно при продвижении времени учитываются ещё не наступившие плановые начала. Для совпадающих времён фиксируется единый порядок обработки по типу события и порядковому номеру; все события текущего момента обрабатываются до проверки новых запусков.

Зарезервированное содержимое означает не только ожидаемые поставки, а сумму физического наличия и ожидаемого выхода уже запущенных операций. Для каждого склада $k$ обозначим через $\mathcal{P}_k(t^{*})$ множество ожидаемых рулонов с индивидуальными идентификаторами. Тогда инвариант имеет вид
\[
\overline{\overline{t}}_k(t^{*})
=\overline{t}_k(t^{*})\mathbin{\dot\cup}\mathcal{P}_k(t^{*}),
\qquad
|\overline{\overline{t}}_k(t^{*})|\leq u^T_k.
\]
Ширина каждого физического или ожидаемого рулона не превышает $r^T_k$. При завершении выпуска рулон переносится из $\mathcal{P}_k$ в $\overline{t}_k$: его ожидаемый статус заменяется фактическим, но запись в $\overline{\overline{t}}_k$ сохраняется и занятая ёмкость не освобождается. Нельзя одновременно учитывать один рулон как фактический и ожидаемый либо удалять его из общей складской занятости при поступлении.

Промежуточные рулоны, получаемые при резке и имеющие нулевой номер конечного заказа, также поступают на склады полуфабриката и подчиняются тем же ограничениям. Поэтому перед запуском резки требуется не только наличие входных тамбуров, но и резервирование всего её промежуточного выхода. Изъятие входов и резервирование выходов проверяются и выполняются как единый переход: место, освобождаемое входами, может использоваться для выходов этой операции. Если допустимого размещения нет, переход не выполняется и входы остаются на складах. Конечная продукция с ненулевыми номерами заказов учитывается в выполнении заказов и не занимает склады полуфабриката $T$; отдельный склад готовой продукции в постановку не входит. Для всех входов и промежуточных выходов сохраняются совместимость по технологическим характеристикам и принадлежность тому же цеху, что и обрабатывающая машина.

Операция может стартовать, только если:
\begin{enumerate}
  \item $t^{*}\geq h(u)$, а все предшествующие операции её машины завершены;
  \item машина свободна и исправна; приостановленная операция удерживает машину до своего завершения;
  \item для резки имеются все требуемые полуфабрикаты по ширине, марке, плотности и числу слоёв;
  \item после изъятия входов и резервирования всех складируемых выходов соблюдены ограничения ширины, ёмкости и помещения каждого склада.
\end{enumerate}
Правила выбора совместимых входов, размещения выходов и разрешения конкуренции одновременно готовых операций фиксируются до эксперимента и совпадают с правилами исходного исполнения. Если исходный план явно задаёт материальные связи, они сохраняются. Изменение этих правил означало бы сравнение разных политик управления, а не только устойчивости одного плана.

Один прогон выполняется следующим образом.
\begin{enumerate}
  \item Установить $t^{*}=H$, пустые склады, исправные свободные машины и незапущенные операции; задать поток случайных чисел. Начальные запасы, если они будут введены как расширение исходных данных, должны одинаково учитываться в исходной и имитационной моделях.
  \item Обработать действительные события текущего момента. Завершение записывает $\overline{h}^{*}(u)$, пополняет фактические склады с заменой ожидаемого статуса, учитывает готовую продукцию и освобождает машину. Отказ приостанавливает операцию, сохраняет остаток обработки и \emph{отменяет прежнее событие завершения}. Для этого можно удалять событие или помечать его номером версии операции; устаревшие события не изменяют состояние. Разыгрывается ремонт и назначается восстановление. Во время ремонта входы не возвращаются, резервы выходов не снимаются.
  \item При восстановлении возобновить остаток обработки, назначить новое завершение и возможный новый отказ. Уже разыгранная длительность всей операции не пересэмплируется. При равенстве наработки до отказа и оставшегося времени приоритет имеет завершение; отказ назначается только при строгом неравенстве $\chi<$ остаток обработки.
  \item В фиксированном порядке запустить доступные операции, выполняя складские переходы атомарно. Для каждой впервые запускаемой операции записать $h^{*}(u)$, разыграть $\widetilde{\pi}(u)$, назначить завершение и возможный отказ. Повторять проверку доступности до отсутствия новых запусков в текущий момент.
  \item Если все операции завершены, закончить прогон с признаком завершения. Иначе найти минимум среди времён действительных будущих событий и всех плановых начал незапущенных операций, \emph{строго больших} $t^{*}$. Если такой момент есть, перейти к нему и повторить обработку. Если будущего момента нет и запуск невозможен, зарегистрировать блокировку и завершить прогон с диагностикой.
  \item Сохранить фактическую историю, признак завершения или блокировки, причины ожидания, состояние складов и показатели прогона.
\end{enumerate}

\begin{center}
\begin{tikzpicture}[
  simstep/.style={draw, rounded corners, align=center, text width=9.5cm, minimum height=9mm},
  simarrow/.style={-{Latex}, thick}, node distance=5mm, font=\small]
  \node[simstep] (init) {Инициализация времени, операций, машин и складов};
  \node[simstep, below=of init] (event) {Обработка действительных событий текущего момента};
  \node[simstep, below=of event] (start) {Возобновление и допустимые запуски; обновление резервов};
  \node[simstep, below=of start] (check) {Все операции завершены: сохранить результат и выйти};
  \node[simstep, below=of check] (next) {Иначе: найти будущее событие или будущее плановое начало};
  \node[simstep, below=of next] (stop) {Будущего момента нет: блокировка, диагностика и выход};
  \draw[simarrow] (init) -- (event);
  \draw[simarrow] (event) -- (start);
  \draw[simarrow] (start) -- (check);
  \draw[simarrow] (check) -- node[right] {не завершены} (next);
  \draw[simarrow] (next) -- node[right] {нет} (stop);
  \draw[simarrow] (next.west) -- ++(-12mm,0) |- node[pos=0.25, left, align=right] {есть:\\сдвиг $t^{*}$} (event.west);
\end{tikzpicture}
\captionof{figure}{Схема дискретно-событийного имитационного прогона}
\label{alg:simulation_run}
\end{center}

Ожидание ещё не наступившего планового начала само по себе не является блокировкой. Не является ею и отсутствие материала, который должен поступить от активной операции: в этом случае существует действительное будущее событие. Блокировка означает отсутствие любого дальнейшего перехода при незавершённом плане и фиксированных правилах исполнения. Диагностика содержит список ожидающих операций, нехватку совместимых полуфабрикатов или складской ёмкости и удерживаемые резервы. Простая временная задержка отлична от структурного взаимного ожидания.

Алгоритм обладает двумя свойствами, обеспечивающими корректность интерпретации результатов. Во-первых, в контрольном нулевом сценарии $\Delta^x=\Delta^y=\nu^x=\nu^y=0$ при одинаковых начальных состояниях, материальных связях и правилах разрешения совпадений имитация в точности воспроизводит исходное расписание: фактические времена операций совпадают с плановыми. Это обеспечивает согласованность имитационной модели с детерминированной постановкой; нулевой сценарий отделён от производственных сценариев. Во-вторых, при фиксированных входных данных и правилах событий прогон полностью детерминирован состоянием генератора~$\zeta$: повторный запуск с тем же значением даёт идентичный результат, обеспечивая воспроизводимость вычислительных экспериментов.

Имитационный эксперимент состоит из $N_{\Theta}$ независимых прогонов; $r$-й прогон использует начальное состояние генератора $\zeta+r-1$. Правило формирования псевдослучайных потоков фиксируется в протоколе эксперимента. Результаты $\Omega_1,\ldots,\Omega_{N_{\Theta}}$ включают фактические расписания и показатели прогонов; состав этих данных и правила агрегации определены ниже. Прогоны, завершившиеся блокировкой, сохраняются с соответствующим признаком.

\textbf{Пример.}
Пусть расписание содержит две операции: производство тамбура на~$m_1$ с плановыми временами $08\!:\!00$--$08\!:\!06$ и его нарезку на~$c_1$ с плановыми временами $08\!:\!06$--$10\!:\!12$; сценарий $\Theta$ задаёт $\Delta^{x}=\Delta^{y}=0{,}1$ без отказов.
Допустим, для операции производства разыграно $\xi=+0{,}05$: фактическая длительность $\widetilde{\pi}=6\cdot1{,}05=6{,}3$~мин, завершение в~$08\!:\!06\!:\!18$.
Операция нарезки не может начаться в плановый момент $08\!:\!06$, так как тамбур ещё не поступил на склад~$\overline{T}$; она запускается в~$08\!:\!06\!:\!18$.
Если для неё разыграно $\xi=-0{,}02$, фактическая длительность составит $126\cdot0{,}98\approx123{,}5$~мин и завершение наступит приблизительно в~$10\!:\!09\!:\!48$ при указанном округлении длительности.
Фактическое расписание~$\widetilde{\overline{\overline{Z}}}$ фиксирует задержку начала нарезки ($18$~с), но общая длительность исполнения в данном прогоне оказывается меньше исходной за счёт отрицательного отклонения длительности нарезки. Следовательно, увеличение симметричного разброса длительностей не гарантирует ухудшения каждого прогона и не задаёт строгого порядка результатов разных сценариев.

Оценим сложность одного прогона. Пусть $N_{\mathrm{ops}}=|\mathcal{U}|$, $F$~--- число отказов в данном прогоне, а $K$~--- число складов. Каждое продвижение времени сопровождается просмотром списка незапущенных операций~$O(N_{\mathrm{ops}})$, учётом складов~$O(K)$ и операциями с очередью~$O(\log(N_{\mathrm{ops}}+F))$; число продвижений времени составляет $O(N_{\mathrm{ops}}+F)$. Итоговая оценка сложности прогона имеет вид
\[
O\!\left((N_{\mathrm{ops}}+F)\left(N_{\mathrm{ops}}+K+
\log(N_{\mathrm{ops}}+F)\right)\right).
\]
Оценка условна по реализованному числу отказов: экспоненциальная модель не задаёт его детерминированной верхней границы. Среднее число отказов связано с суммарной наработкой всех машин с учётом их интенсивностей, а не только с календарной длительностью плана. Стоимость эксперимента из $N_{\Theta}$ прогонов при общей оценке числа отказов составляет
\[
O\!\left(N_{\Theta}(N_{\mathrm{ops}}+F)\left(N_{\mathrm{ops}}+K+
\log(N_{\mathrm{ops}}+F)\right)\right),
\]
то есть в $N_{\Theta}$ раз больше стоимости одного прогона; при различающихся числах отказов стоимости отдельных прогонов суммируются. Имитация выполняется для уже построенного плана и не входит во внутренний цикл оптимизации.

Для завершившегося прогона $r$ определим длительность исполнения относительно начала производства и её исходный аналог:
\[
\Lambda_r=\max_{u\in\mathcal{U}}\overline{h}^{*}_r(u)-H,
\qquad
\Lambda_0=\max_{u\in\mathcal{U}}\overline{h}(u)-H.
\]
Таким образом, календарный момент завершения равен $H+\Lambda_r$, а $\Lambda_r$ измеряется в единицах времени, а не является датой. Подписанное отклонение и неотрицательная задержка относительно исходной продолжительности равны
\[
\delta\Lambda_r=\Lambda_r-\Lambda_0,\qquad
D^{\mathrm{plan}}_r=\max\{0,\Lambda_r-\Lambda_0\}.
\]
Отрицательное $\delta\Lambda_r$ означает более раннее завершение; оно не компенсирует опоздание отдельного заказа при подсчёте его запаздывания.

Для заказа $i$ обозначим через $\mathcal{U}_i$ операции получения конечных рулонов этого заказа в полном количестве $q_i$. В завершённом прогоне
\[
C_{i,r}=\max_{u\in\mathcal{U}_i}\overline{h}^{*}_r(u),\qquad
D_{i,r}=\max\{0,C_{i,r}-h_i\},
\]
\[
Q^{\mathrm{late}}_r=\sum_{i=1}^{Q}\mathbf{1}\{D_{i,r}>0\},\qquad
D_{\Sigma,r}=\sum_{i=1}^{Q}D_{i,r},\qquad
D_{\max,r}=\max_{1\leq i\leq Q}D_{i,r}.
\]
Здесь $\mathbf{1}\{\cdot\}$~--- индикатор истинности условия. $Q^{\mathrm{late}}_r$ считает заказы, а не отдельные рулоны; $D_{\Sigma,r}$ характеризует совокупное запаздывание без взвешивания объёмами, $D_{\max,r}$~--- наибольшее запаздывание одного заказа. Для незавершённого заказа момент $C_{i,r}$ ещё не наблюдался: его нельзя заменять временем остановки прогона и объявлять заказ выполненным вовремя.

Пусть $c_r$~--- индикатор полного завершения прогона, $b_r$~--- индикатор блокировки. Для серии из $N_{\Theta}$ прогонов положим $\mathcal{C}_{\Theta}=\{r:c_r=1\}$ и $n_c=|\mathcal{C}_{\Theta}|$. Определим
\[
\gamma_{\mathrm{ok}}(G,\Theta)=
\frac{1}{N_{\Theta}}\sum_{r\in\mathcal{C}_{\Theta}}
\mathbf{1}\{Q^{\mathrm{late}}_r=0\},\qquad
\gamma_{\mathrm{dl}}(G,\Theta)=\frac{1}{N_{\Theta}}\sum_{r=1}^{N_{\Theta}}b_r.
\]
Первая величина~--- эмпирическая доля всех прогонов, в которых план полностью завершён без опозданий, вторая~--- эмпирическая доля блокировок. Они не дополняют друг друга до единицы: возможны завершённые прогоны с опозданиями. Это выборочные оценки при заданном сценарии, а не гарантированные вероятности реального исполнения.

Для каждого $X_r\in\{\Lambda_r,\delta\Lambda_r,D^{\mathrm{plan}}_r,Q^{\mathrm{late}}_r,D_{\Sigma,r},D_{\max,r}\}$ среднее определяется только по завершённым прогонам:
\[
\overline{X}_{\mathrm{comp}}=\frac{1}{n_c}\sum_{r\in\mathcal{C}_{\Theta}}X_r,
\qquad n_c>0.
\]
Для каждого заказа по завершённым прогонам дополнительно оценивается вероятность нарушения его срока:
\[
\widehat{p}^{\mathrm{late}}_i(G,\Theta\mid c=1)
=\frac{1}{n_c}\sum_{r\in\mathcal{C}_{\Theta}}\mathbf{1}\{D_{i,r}>0\},
\qquad n_c>0.
\]
Это условная эмпирическая вероятность опоздания заказа; вместе с ней приводятся среднее и максимальное запаздывание данного заказа по той же выборке. Незавершённые заказы не засчитываются как своевременные.

По той же условной выборке приводятся стандартное отклонение (при $n_c\geq2$), минимум, максимум, медиана и квантили уровней $0{,}05$ и $0{,}95$. Если $n_c=0$, эти показатели не определены, а не равны нулю. Вместе с ними обязательно указываются $n_c$, $N_{\Theta}$ и $\gamma_{\mathrm{dl}}$: малое условное среднее при частых блокировках не означает высокое качество плана. При необходимости после завершённого прогона можно анализировать фактическую цепочку критических ожиданий, отделяя её от исходных зависимостей расписания.

В описанной процедуре каждый прогон заканчивается завершением или блокировкой. Если при проведении эксперимента дополнительно устанавливается предел модельного времени либо вычислительного бюджета, остановленные по нему прогоны учитываются отдельно как незавершённые цензурированные наблюдения. Они не включаются в условные средние, не считаются успешными и не объявляются блокировками без выполнения её критерия. Знаменатели $\gamma_{\mathrm{ok}}$ и $\gamma_{\mathrm{dl}}$ по-прежнему равны числу всех начатых прогонов, а доля остановок по лимиту публикуется отдельно.

Показатели оборудования и складов вычисляются по временной истории. Пусть $\tau_r$~--- момент окончания наблюдения (завершения, блокировки или остановки по лимиту), $W_r=\tau_r-H>0$. Для машины $a$ обозначим суммарное время выполнения операций через $B_{a,r}$, ремонта~--- через $R_{a,r}$, а ожидания~--- через $I_{a,r}$. Состояния взаимно исключаются, поэтому
\[
B_{a,r}+R_{a,r}+I_{a,r}=W_r,\qquad
\eta_{a,r}=\frac{B_{a,r}}{W_r},\quad
\eta^{\mathrm{rep}}_{a,r}=\frac{R_{a,r}}{W_r},\quad
\eta^{\mathrm{idle}}_{a,r}=\frac{I_{a,r}}{W_r}.
\]
$B_{a,r}$ включает предусмотренные длительностью операции технологические переходы; без дополнительной детализации его нельзя называть временем чистой резки или чистого выпуска. Время ожидания можно разложить по причинам (плановый старт, материал, ёмкость) при фиксированном правиле приоритета причин, исключающем двойной счёт. Дополнительно фиксируются число отказов и суммарная длительность ремонта каждой машины.

Пусть $n_{k,r}(t)=|\overline{t}_k(t)|$~--- физическое наличие, $v_{k,r}(t)=|\overline{\overline{t}}_k(t)|$~--- наличие с ожидаемыми поступлениями на складе $k$. При $u^T_k>0$ определим
\[
\overline{\rho}_{k,r}=\frac{1}{u^T_k W_r}\int_H^{\tau_r}n_{k,r}(t)\,dt,
\qquad
\overline{\rho}^{\mathrm{res}}_{k,r}=\frac{1}{u^T_k W_r}\int_H^{\tau_r}v_{k,r}(t)\,dt,
\]
\[
\rho^{\max}_{k,r}=\max_{H\leq t\leq\tau_r}\frac{n_{k,r}(t)}{u^T_k},\qquad
\rho^{\mathrm{res},\max}_{k,r}=\max_{H\leq t\leq\tau_r}\frac{v_{k,r}(t)}{u^T_k}.
\]
Это соответственно средняя и пиковая доли физически и с учётом резервов занятой ёмкости. Доли времени полной загрузки склада определяются отдельно для физического наличия и занятости с резервами:
\[
\phi^{\mathrm{full}}_{k,r}=\frac{1}{W_r}\int_H^{\tau_r}\mathbf{1}\{n_{k,r}(t)=u^T_k\}\,dt,
\qquad
\phi^{\mathrm{res,full}}_{k,r}=\frac{1}{W_r}\int_H^{\tau_r}\mathbf{1}\{v_{k,r}(t)=u^T_k\}\,dt.
\]
При нулевой ёмкости нормированные показатели не определяются; контролируется отсутствие поступлений. Интегралы вычисляются по кусочно-постоянным состояниям между событиями. Для $W_r=0$ временные доли также не определяются. При завершении $W_r=\Lambda_r$; показатели короткого заблокированного прогона характеризуют только наблюдённый фрагмент. Поэтому для сравнения планов приводятся условные агрегаты загрузки по завершённым прогонам, а незавершённые истории и их длительности рассматриваются отдельно, без преимущества за раннюю остановку.

Таким образом, данные отдельного прогона имеют вид
\[
\Omega_r=\bigl(\widetilde{\overline{\overline{Z}}}_r,c_r,b_r,
\Lambda_r,Q^{\mathrm{late}}_r,D_{\Sigma,r},D_{\max,r},\mathcal{H}_r\bigr),
\]
где $\widetilde{\overline{\overline{Z}}}_r$~--- фактическое расписание с зарегистрированными началами и окончаниями операций, а $\mathcal{H}_r$ объединяет временную историю событий, причины ожиданий, диагностику остановки и определённые выше показатели оборудования, отдельных заказов и складов. Для незавершённого прогона сохраняется наблюдённая часть расписания; величины, требующие полного завершения, отмечаются как неопределённые, а не заменяются нулями. Совокупность $\Omega_1,\ldots,\Omega_{N_{\Theta}}$ служит исходными данными для статистических оценок эксперимента.

Неопределённость оценок описывается доверительными интервалами с заранее выбранным уровнем, например $0{,}95$. Для долей применимы биномиальные интервалы; отсутствие наблюдённых блокировок не означает нулевой верхней границы их вероятности. Для условных средних и квантилей можно использовать повторные выборки завершённых прогонов, явно указывая условность оценки и объём $n_c$. При малом $n_c$ оценка верхнего квантиля нестабильна; увеличение числа повторений уменьшает выборочную неопределённость, но не устраняет ошибку выбранной модели случайных факторов.

Для сравнительной оценки фиксируются три производственных сценария и~расширяемые параметры модели. До запуска сравнительных опытов фиксируются три отдельных набора параметров:
\[
\Theta_s=(\Delta^x_s,\Delta^y_s,\nu^x_s,\nu^y_s,\theta_s),
\qquad s\in\{\mathrm{opt},\mathrm{real},\mathrm{pess}\}.
\]
В пределах каждого сценария набор не меняется между планами и повторениями; случайны реализации, а не сами сценарные параметры. Значения, единицы времени, источники оценок, число повторений и способ формирования потоков случайных чисел должны быть заданы в протоколе эксперимента до получения результатов.
\begin{enumerate}
  \item \emph{Оптимистичный сценарий} $\Theta_{\mathrm{opt}}$ описывает благоприятные, но не идеально детерминированные условия: малый разброс длительностей, редкие отказы и быстрое восстановление. Хотя бы один параметр возмущения отличен от нуля; этот сценарий не заменяет контрольный нулевой прогон.
  \item \emph{Реалистичный сценарий} $\Theta_{\mathrm{real}}$ предназначен для типичных условий конкретного производства. Его параметры калибруются по данным о фактических длительностях, наработках до отказа и ремонтах, с проверкой пригодности выбранных распределений. Калибровка связывает сценарные предположения с типичными условиями предприятия; источники данных и метод оценки параметров фиксируются в протоколе эксперимента.
  \item \emph{Пессимистичный сценарий} $\Theta_{\mathrm{pess}}$ задаёт стрессовые предположения: более широкий разброс длительностей, более частые отказы и более длительное восстановление. Он служит проверкой чувствительности плана, а не доказанной границей наихудшего возможного исполнения.
\end{enumerate}
Для согласованной интерпретации задаются порядковые соотношения
\[
\Delta^e_{\mathrm{opt}}\leq\Delta^e_{\mathrm{real}}\leq\Delta^e_{\mathrm{pess}},\quad
\nu^e_{\mathrm{opt}}\leq\nu^e_{\mathrm{real}}\leq\nu^e_{\mathrm{pess}},\quad e\in\{x,y\},
\qquad
\theta_{\mathrm{opt}}\leq\theta_{\mathrm{real}}\leq\theta_{\mathrm{pess}}.
\]
Наборы должны различаться хотя бы по одному параметру. Симметричное увеличение $\Delta^e$ сохраняет среднюю длительность работы $\pi(u)$, увеличивая разброс, поэтому названия сценариев характеризуют предположения, но не гарантируют монотонности всех наблюдаемых метрик. Среднее, медиана и квантили вычисляются \emph{внутри каждого} сценария; они не являются определениями оптимистичного, реалистичного и пессимистичного сценариев.

Параметризация допускает дальнейшее расширение модели. В качестве направлений расширения можно рассмотреть: систематические относительные сдвиги $b^{(e)}$ в формуле $\widetilde{\pi}(u)=\pi(u)(1+b^{(e)}+\xi_u)$ при $1+b^{(e)}-\Delta^{(e)}>0$; индивидуальные параметры разброса, отказов и восстановления для каждой машины; зависимость от марки и вида операции; отдельный календарь плановых ремонтов и переналадок. Для календарных расширений требуется определить допустимость прерываний и исключить повторный учёт переналадок, уже включённых в $\pi(u)$. Такие изменения требуют дополнительных исходных данных, проверки согласованности событий и самостоятельной валидации; они не следуют автоматически из базовой имитационной модели.

\section{Сравнение кандидатов и выбор плана}\label{sec:model/selection}
После завершения численного поиска его результатом для имитационной оценки служит набор
\[
\mathcal{A}=\{G^{(1)},\ldots,G^{(L)}\},\qquad 1\leq L\leq A_{\max},
\]
где $A_{\max}$~--- желаемое максимальное число различных кандидатов, а $L$~--- фактически найденное число. Это обозначение не следует смешивать с максимальным числом слоёв $L_{\max}$ в оценке сложности исходной целевой функции. При отсутствии допустимых планов набор пуст и этап выбора не выполняется. Кандидаты различаются содержанием глобального плана, а не только идентификатором; равенство значений $\mu$ не означает равенства самих планов.

Набор формируется ограниченным архивом различных лучших по исходному лексикографическому значению $\mu$ планов, встреченных в ходе поиска. В эволюционном методе архив пополняется при рассмотрении планов до и после мутации и отбора, чтобы перспективный вариант не исчезал только из-за изменения текущей популяции. При переполнении сохраняются не более $A_{\max}$ исходно лучших различных кандидатов с фиксированным правилом разрешения равенства оценок. Имитация выполняется только после завершения поиска и не участвует в операторах мутации или отбора. Подробные численные процедуры относятся к следующей главе; здесь задано требование к их выходу для последующего сравнения устойчивости.

Такой архив не является множеством всех допустимых планов или гарантированным фронтом устойчивых решений. Исходно менее выгодный план, не попавший в архив, может оказаться устойчивее сохранённых. Поэтому окончательный выбор относится только к $\mathcal{A}$ и заданным сценариям и не гарантирует глобальной робастности.

Правило сравнения, сценарные параметры и число повторений устанавливаются до опытов. Например, задаются минимально приемлемая доля своевременного завершения $p_{\min}\in[0,1]$ и максимально приемлемая доля блокировок $b_{\max}\in[0,1]$. Для каждого сценария строятся нижняя доверительная граница $\underline{\gamma}_{\mathrm{ok}}$ и верхняя граница $\overline{\gamma}_{\mathrm{dl}}$. Кандидат проходит статистический фильтр, если для всех трёх сценариев
\[
\underline{\gamma}_{\mathrm{ok}}(G,\Theta_s)\geq p_{\min},\qquad
\overline{\gamma}_{\mathrm{dl}}(G,\Theta_s)\leq b_{\max}.
\]
Это консервативное правило принятия решения по данным, а не гарантия выполнения порогов истинными вероятностями. Уровень интервалов и поправка на множественное сравнение планов и сценариев, если она применяется, также задаются заранее.

Среди прошедших фильтр вариантов можно минимизировать наихудший по сценариям условный квантиль суммарного запаздывания:
\[
R_{95}(G)=\max_{s\in\{\mathrm{opt},\mathrm{real},\mathrm{pess}\}}
\widehat{q}_{0.95}\!\left(D_{\Sigma}\mid c=1,\ G,\Theta_s\right).
\]
Здесь $\widehat{q}_{0.95}$~--- выборочная оценка $95$-го процентиля по завершённым прогонам соответствующего сценария. При отсутствии завершённых прогонов квантиль не определён и кандидат не может быть выбран по этому правилу. Для близких оценок учитываются интервалы неопределённости квантилей; заранее устанавливаются минимальный объём завершённых прогонов, правило дополнительного моделирования и правило разрешения статистически неразличимых результатов, например по исходной $\mu$, затем по фиксированному порядку планов. Если ни один кандидат не проходит фильтр, сообщается об отсутствии подтверждённо приемлемого варианта в оценённом наборе, а не о доказанной невозможности устойчивого планирования вообще.

Таким образом, лексикографическое сравнение без весов остаётся правилом \emph{исходного поиска}. Сценарные пороги и квантиль запаздывания образуют отдельное правило \emph{последующего выбора} по устойчивости и не изменяют математическую целевую функцию внутри численного метода. Для ограничения эффекта выбора по случайно удачной выборке окончательную оценку выбранного плана целесообразно подтверждать независимой серией прогонов по заранее заданному протоколу.

\section{Выводы по главе}\label{sec:model/statement_conclusion}
В данной главе сформирована схема оценки устойчивости производственных планов для многомашинной системы целлюлозно-бумажного производства: от получения нескольких исходно допустимых альтернатив до их сценарного сравнения и выбора. Математическая модель задачи раскроя и формирования расписания задаёт основу этой схемы, определяя состав операций, материальные связи и ограничения исполнения.

Основные результаты главы:
\begin{enumerate}
    \item Введены входные параметры, описывающие заказы $O$ (включая многослойные заказы и~индикатор БРС~$g_i$), оборудование ($M$, $C$, $T$), производственные помещения ($S$ с~параметрами принадлежности ${\rho^m}_i$, ${\rho^c}_j$, ${\rho^T}_k$) и~технологические константы ($B$, $V$, $H$).
    Определены переменные решения~--- план производства $\overline{\overline{X}}$, план нарезки $\overline{\overline{Y}}$ и расписание $\overline{\overline{Z}}$, объединённые в глобальный план $G$.
    Трёхкомпонентная структура~$G$ обеспечивает совместный учёт выпуска, раскроя и~календарного планирования в~рамках единой модели.
    \item Сформулирована векторная целевая функция $\mu(\overline{\overline{Y}}) = \langle \omega(\overline{\overline{Y}}),\, \varphi(\overline{\overline{Y}}),\, \tau(\overline{\overline{Y}}) \rangle \to \min$, отражающая три ключевых критерия качества плана: минимизацию отходов, сбалансированность загрузки станков и сокращение переналадок. Обоснован выбор лексикографического упорядочения критериев.
    \item Определена система из десяти ограничений, обеспечивающих технологическую корректность плана: соблюдение порядка марок, допустимость раскроев по ширине, полноту расписания, выполнение сроков заказов, наличие полуфабриката для нарезки, непереполнение складов и размещение взаимодействующего оборудования в одном производственном помещении (цеховое ограничение).
    Ограничения по~складам (пп.~7--9) учитывают как содержимое ($\overline{T}$), так и~зарезервированное содержимое ($\overline{\overline{T}}$) складов, что отличает модель от~постановок с~агрегированными переменными запасов \cite{nonas2000combined,leao2016decomposition}.
    \item Показано, что рассматриваемая задача относится к классу NP-трудных, а её размерность определяется количеством заказов $Q$, числом единиц оборудования $I$, $J$ и длиной планов $|\overline{X}_i|$, $|\overline{Y}_i|$.
    \item Разработана система дискретно-событийного имитационного моделирования исполнения фиксированного плана при случайном разбросе длительностей, отказах и ремонтах, позволяющий оценивать устойчивость построенного плана. Определены условия запуска, приостановки и возобновления операций, сохранения складских резервов и обнаружения блокировок; ожидание будущего планового начала отделено от невозможности дальнейшего исполнения. В контрольном нулевом сценарии алгоритм воспроизводит исходное расписание; фиксация состояния генератора случайных чисел обеспечивает воспроизводимость экспериментов.
    \item Введены показатели длительности относительно начала производства, задержки относительно исходной продолжительности, числа просроченных заказов, суммарного и максимального запаздывания, а также доли своевременных завершений и блокировок. Заданы временные знаменатели показателей машин и складов; условные оценки по завершённым прогонам отделены от статистики незавершения по всей серии.
    \item Разграничены оптимистичный, реалистичный и пессимистичный сценарии как три фиксируемых набора параметров, а также отдельный контрольный нулевой сценарий. Для реалистичного сценария указана необходимость калибровки по производственным данным; индивидуальные параметры машин, сдвиги длительностей и календарные воздействия сформулированы как направления расширения.
    \item Задано сравнение ограниченного набора различных исходных кандидатов после завершения поиска: по заранее установленным требованиям к своевременному завершению и блокировкам, затем по квантилю запаздывания с учётом неопределённости оценок. Выбор ограничен рассмотренными планами и сценарными предположениями и не означает гарантии глобальной устойчивости.
\end{enumerate}

Научная новизна модели определяется совокупностью одновременно учитываемых аспектов, не~встречающейся в~известных работах: несколько ПРС и~БРС в~одной модели, многослойные заказы ($l_i>1$), двухэтапная нарезка малых форматов через БРС ($g_i=1$), динамика складских запасов на~уровне отдельных тамбуров, цеховое ограничение на~размещение взаимодействующего оборудования в~одном помещении и~построение согласованного расписания~$\overline{\overline{Z}}$ с~контролем состояния складов.

Модель согласована с численными методами, программной реализацией и вычислительными экспериментами. Построенная исходная модель и процедура имитационной оценки используются в последующих главах как основа для разработки численных методов, их программной реализации и постановки вычислительных экспериментов. Численные методы формируют набор альтернатив по исходным критериям, а их устойчивость оценивается отдельно после завершения поиска.

\FloatBarrier
